\subsubsection*{Electric and Magnetic Fields}
The next step is to spell out the electromagnetic field content that was previously presented in compact STA form. These identities make the relation between the vector fields and the single field paravector fully explicit. See Eqs.\ \eqref{eq:dalembert}, \eqref{eq:F-Phi}, and \eqref{eq:maxwell}. Here \(V,\rho,S,q\in\R\), \(\mathbf A,\mathbf E,\mathbf B,\mathbf J\in\R^3\), \(\Phi=V+\mathbf A\cdot\sigv\in\R\oplus\R^3\), and \(F=(\mathbf E+i\mathbf B)\cdot\sigv\in\{0\}\oplus\C^3\).
\noindent\textit{Introduce the following vector notation.}\par
\begin{align}
\delr{} &= \{\partial /\partial r_{1}, \partial /\partial r_{2}, \partial /\partial r_{3}\},\notag \\
\partial _{i} &= \partial / \partial r_{i},\notag \\
\delr{}^{2} &= \partial _{1}^{2} + \partial _{2}^{2} + \partial _{3}^{2},\notag \\
\partial &= \partial _{t} + \sigv \cdot \delr{},\notag \\
&= \partial _{t} + \sigma _{1}\partial _{1}+ \sigma _{2}\partial _{2} + \sigma _{3}\partial _{3}.
\end{align}
\noindent\textit{Introduce the following vectors.}\par
\begin{align}
\mathbf E &= \{\mathbf E_{1}, \mathbf E_{2}, \mathbf E_{3}\}\qquad \text{electric field},\notag \\
\mathbf B &= \{\mathbf B_{1}, \mathbf B_{2}, \mathbf B_{3}\}\qquad \text{magnetic field},\notag \\
\mathcal{F} &= \mathbf E + i \mathbf B\qquad \text{electromagnetic field},\notag \\
F &= (\mathbf E + i \mathbf B) \cdot \sigv , i = \sigma _{1}\sigma _{2}\sigma _{3},\notag \\
&= \mathcal{F} \cdot \sigv.
\end{align}
\begin{align}
F^{*} &= (-\mathbf E + i \mathbf B) \cdot \sigv,\notag \\
&= -\mathcal{F}^{*} \cdot \sigv.
\end{align}

\noindent\textit{Introduce the source-density paravector.}\par
\begin{align}
D &= \rho + \mathbf J \cdot \sigv, \\
D^{*} &= \rho - \mathbf J \cdot \sigv.
\end{align}
Let \(\rho\) denote the scalar charge density.\par
Let \(\mathbf J\) denote the current-density three-vector.\par
\begin{align}
\partial F &= D^{*}\qquad \text{Maxwell's equations},\notag \\
&= \rho - \mathbf J \cdot \sigv.
\end{align}
\begin{equation}
\partial ^{*} F^{*} = D.
\end{equation}
\noindent\textit{Expanding the left-hand side gives.}\par
\begin{align}
\partial F &= (\partial _{t} + \sigv \cdot \delr{}) ( (\mathbf E + i \mathbf B) \cdot \sigv ),\notag \\
&= ( \delr{} \cdot \mathbf E ),\notag \\
&\quad + ( \delr{} \cdot \mathbf B ) i,\notag \\
&\quad + (\partial _{t} \mathbf E - \delr{} \times \mathbf B) \cdot \sigv,\notag \\
&\quad + (\partial _{t} \mathbf B + \delr{} \times \mathbf E) \cdot (i \sigv ).
\end{align}
\noindent\textit{Rearranging the vector equations gives.}\par
\begin{align}
\delr{} \times \mathbf B &= \mathbf J + \partial _{t} \mathbf E, \\
\delr{} \times \mathbf E &= -\partial _{t} \mathbf B.
\end{align}

\subsubsection*{Lorentz Transformation of the Field}
Once the electromagnetic field has been assembled into the single paravector \(F\), the boost law from Eq.\ \eqref{eq:field-transform} can be expanded directly in components. For a boost
\(L=\exp(\tfrac12\theta\,\mathbf u\cdot\sigv)\) with \(\mathbf u^{2}=1\), \(\beta=\tanh(\theta)\), and \(\gamma=\cosh(\theta)\), the parallel and perpendicular field parts relative to \(\mathbf u\) obey
\begin{align}
F' &= L^{*} F L, \\
\eqtext{par}( \mathbf E ' ) &= \eqtext{par}( \mathbf E ), \\
\eqtext{par}( \mathbf B ' ) &= \eqtext{par}( \mathbf B ), \\
\eqtext{prp}( \mathbf E ' ) &= \gamma ( \eqtext{prp}( \mathbf E ) + \beta \mathbf u \times \mathbf B ), \\
\eqtext{prp}( \mathbf B ') &= \gamma ( \eqtext{prp}( \mathbf B ) - \beta \mathbf u \times \mathbf E ).
\end{align}
\noindent\textit{These identities can be compressed into the single relation.}\par
\begin{align}
\mathbf E ' &= \eqtext{par}(\mathbf E) + \gamma ( \eqtext{prp}(\mathbf E) + \beta \mathbf u \times \mathbf B ), \\
\mathbf B ' &= \eqtext{par}(\mathbf B) + \gamma ( \eqtext{prp}(\mathbf B) - \beta \mathbf u \times \mathbf E ), \\
\partial ' &= L^{*} \partial L, \\
\partial ' F' &= L^{*} (\rho - \mathbf J \cdot \sigv ) L.
\end{align}

\subsubsection*{Wave Operator and d’Alembertian}
The wave operator is written out here in full so that its factorization into the d'Alembertian and the component field equations can be read directly. See Eq.\ \eqref{eq:dalembert}.
\begin{align}
\partial &= \partial _{t} + \sigv \cdot \delr{}, \\
\partial ^{*} &= \partial _{t} - \sigv \cdot \delr{}.
\end{align}
\noindent\textit{The d'Alembertian is.}\par
\begin{align}
\square &:= \partial _{t}^{2} - \delr{}^{2},\notag \\
&= \partial \partial ^{*} = \partial ^{*} \partial,\notag \\
&= \langle\!\langle \partial \rangle\!\rangle^{2}.
\end{align}
\noindent\emph{Note.} This expansion uses the sigma-algebra conventions fixed earlier in the manuscript.\par
\begin{equation}
\partial ^{2} = \partial _{t}^{2} + \delr{}^{2} + 2 \partial _{t} \delr{} \cdot \sigv.
\end{equation}
\noindent\textit{This differs from }\(\square\)\textit{.}\par
\noindent\emph{Note.} Sign conventions for \(F\) vary across the literature; the present manuscript keeps the convention fixed by Eqs.\ \eqref{eq:F-Phi} and \eqref{eq:maxwell}.\par
\noindent\emph{Note.} Some authors instead use the metric signature \((- ,+,+,+)\).\par
\noindent\emph{Note.} Some authors denote the wave operator with a different box-symbol convention.\par
\noindent\textit{The corresponding wave equation is.}\par
\begin{align}
\square F &= \partial ^{*} \partial F,\notag \\
&= \partial ^{*} (\rho - \mathbf J \cdot \sigv ).
\end{align}
\begin{align}
\square F^{*} &= \partial (\rho + \mathbf J \cdot \sigv ),\notag \\
( \square \mathbf E + i \square \mathbf B ) \cdot \sigv,\notag \\
&= (\partial _{t} \rho + \delr{} \cdot \mathbf J) + (-\partial _{t} \mathbf J - \delr{} \rho + i (\delr{} \times \mathbf J)) \cdot \sigv.
\end{align}
\begin{equation}
\partial _{t} \rho = - \delr{} \cdot \mathbf J\qquad \text{conservation of charge}.
\end{equation}
If \(\rho = 0\) and \(\mathbf J = 0\), then the following relation holds.\par
\begin{align}
\square \mathbf E &= 0 \Rightarrow \partial _{t}^{2} \mathbf E = \delr{}^{2} \mathbf E, \\
\square \mathbf B &= 0 \Rightarrow \partial _{t}^{2} \mathbf B = \delr{}^{2} \mathbf B.
\end{align}

\subsubsection*{Potential Fields}
The potential-based formulation is expanded here because it is the point at which gauge freedom, source reconstruction, and the emergence of the electric and magnetic fields all become transparent at once. See Eqs.\ \eqref{eq:potential-gradient} and \eqref{eq:E-B-from-potential}. Here \(V,S,\rho,\lambda\in\R\), \(\mathbf A,\mathbf E,\mathbf B,\mathbf J\in\R^3\), \(\Phi\in\R\oplus\R^3\), and \(F\in\{0\}\oplus\C^3\subset\C\oplus\C^3\).
\begin{align}
F^{*} &= (-\mathbf E + i \mathbf B) \cdot \sigv,\notag \\
\Phi &= V + \mathbf A \cdot \sigv,\notag \\
\partial \Phi &= \partial _{t} V + \delr{} \cdot \mathbf A,\notag \\
&\quad + ( \partial _{t} \mathbf A + \delr{} V ) \cdot \sigv,\notag \\
&\quad + i ( \delr{} \times \mathbf A ) \cdot \sigv.
\end{align}
\begin{align}
S &:= \eqtext{scalar}( \partial \Phi ),\notag \\
&= \partial _{t} V + \delr{} \cdot \mathbf A.
\end{align}
\noindent\textit{The target component identities are.}\par
\begin{align}
\mathbf E &= - \partial _{t} \mathbf A - \delr{} V, \\
\mathbf B &= \delr{} \times \mathbf A.
\end{align}
\noindent\textit{Consider the following identity.}\par
\begin{align}
\partial \Phi &= S + F^{*}, \\
F^{*} &= \partial \Phi - S, \\
\eqtext{scalar}(F^{*}) &= \eqtext{scalar}(\partial \Phi - S) = 0, \\
\eqtext{vector}(F^{*}) &= \eqtext{vector}( \partial \Phi ).
\end{align}
\noindent\emph{Note.} always differentiate first.\par
\begin{equation}
(\partial \Phi )^{*\mathrm{R}} \neq \Phi ^{*\mathrm{R}} \partial ^{*\mathrm{R}}.
\end{equation}
\noindent\textit{Equivalently, without introducing \(S\), one may write.}\par
\begin{equation}
F^{*} = (\partial \Phi - (\partial \Phi )^{*\mathrm{R}}) / 2.
\end{equation}

\subsubsection*{Sources Derived from the Potential}
\noindent\textit{Recall the previously established definitions.}\par
\begin{align}
F &= (\mathbf E + i \mathbf B) \cdot \sigv,\notag \\
\Phi &= V + \mathbf A \cdot \sigv,\notag \\
S &= \partial _{t} V + \delr{} \cdot \mathbf A,\notag \\
F^{*} &= \partial \Phi - S \Rightarrow F = \partial ^{*} \Phi ^{*} - S,\notag \\
\partial F &= \partial \partial ^{*} \Phi ^{*} - \partial S,\notag \\
&= \square \Phi ^{*} - \partial S,\notag \\
&= ( \square V - \partial _{t} S) + (- \square \mathbf A - \delr{} S) \cdot \sigv,\notag \\
&= \rho - \mathbf J \cdot \sigv.
\end{align}
\begin{align}
\rho &= \square V - \partial _{t} S, \\
\mathbf J &= \square \mathbf A + \delr{} S.
\end{align}
If \(S = 0,\), then the following relation holds.\par
\noindent\textit{Lorenz gauge condition.}\par
\begin{align}
\partial _{t} V &= -\delr{} \cdot \mathbf A, \\
\rho &= \square V, \\
\mathbf J &= \square \mathbf A.
\end{align}

\subsubsection*{Gauge Transform of the Potential}
\begin{align}
\eqtext{scalar}( \partial \Phi ) &= S, \\
S &= \partial _{t} V + \delr{} \cdot \mathbf A.
\end{align}
Let \(\lambda\) denote scalar field.\par
\begin{align}
\Phi ' &= \Phi + \partial ^{*} \lambda,\notag \\
\partial \Phi ' &= \partial (\Phi + \partial ^{*}\lambda ),\notag \\
&= \partial \Phi + \partial (\partial ^{*}\lambda ),\notag \\
&= \partial \Phi + \square \lambda.
\end{align}
\noindent\emph{Note.} The quantity \(\square\lambda\) is again a scalar field.\par
\begin{align}
\eqtext{scalar}( \partial \Phi ' ) &= S + \square \lambda, \\
\eqtext{vector}( \partial \Phi ' ) &= \eqtext{vector}( \partial \Phi ), \\
F' &= F = (\mathbf E + i \mathbf B) \cdot \sigv.
\end{align}

\subsubsection*{Proper Force and Lorentz Dynamics}
\begin{align}
m U &= m dX/ds,\notag \\
&= m \gamma + m \gamma \mathbf v \cdot \sigv.
\end{align}
\begin{equation}
d/ds = \gamma d/dt.
\end{equation}
\noindent\textit{The time component gives the relativistic work law.}\par
\begin{equation}
d/dt (m \gamma ) = q \mathbf v \cdot \mathbf E.
\end{equation}
\noindent\emph{Note.} Here \(\mathbf E\) denotes the electric field, not the energy operator \(E\).\par
\noindent\emph{Note.} The magnetic field contributes no work term in the scalar component.\par
\noindent\textit{The spatial component gives the relativistic Lorentz force law.}\par
\begin{align}
\mathbf P &:= m \gamma \mathbf v, \\
d\mathbf P/dt &= q \mathbf E + q \mathbf v \times \mathbf B.
\end{align}
\noindent\textit{These terms combine into the single relation.}\par
\begin{equation}
d/dt (m U) = q \mathbf v \cdot \mathbf E + (q \mathbf E + q \mathbf v \times \mathbf B) \cdot \sigv.
\end{equation}
\noindent\textit{Consider the following identity.}\par
\begin{align}
F U / \gamma &= ((\mathbf E + i \mathbf B) \cdot \sigv ) ( 1 + \mathbf v \cdot \sigv ),\notag \\
&= \mathbf E \cdot \sigv + i \mathbf B \cdot \sigv + (\mathbf E \cdot \sigv ) (\mathbf v \cdot \sigv ),\notag \\
&\quad + i (\mathbf B \cdot \sigv )(\mathbf v \cdot \sigv ),\notag \\
&= \mathbf E \cdot \sigv + i \mathbf B \cdot \sigv + \mathbf E \cdot \mathbf v,\notag \\
&\quad + i (\mathbf E \times \mathbf v) \cdot \sigv + i \mathbf B \cdot \mathbf v - (\mathbf B \times \mathbf v) \cdot \sigv.
\end{align}
\begin{align}
\eqtext{real}( F U ) &= \tfrac{1}{2} ( F U + U F^{\mathrm{R}} ),\notag \\
&= \gamma (\mathbf v \cdot \mathbf E + ( \mathbf E + \mathbf v \times \mathbf B ) \cdot \sigv ).
\end{align}
\noindent\emph{Note.} The imaginary part of \(FU\) exchanges the electric and magnetic roles and is therefore not the sector used in the physical force law.\par
\begin{align}
m dU/ds &= q \eqtext{real}( F U ),\notag \\
&= \tfrac{1}{2} q ( F U + U F^{\mathrm{R}} ).
\end{align}
\begin{equation}
m ddX/ds^{2} = q \eqtext{real}( F dX/ds )\qquad \text{4-force}.
\end{equation}

\subsubsection*{Variational Derivation of the Lorentz Force}
\noindent\textit{Recall the previously established definitions.}\par
\begin{align}
\Phi &= V + \mathbf A \cdot \sigv,\notag \\
U &= dX/ds,\notag \\
\gamma &= dt/ds,\notag \\
\text{action} &= -m \int ds - q \int \eqtext{scalar}( \Phi ^{*} U ) ds,\notag \\
&= \int ( -m / \gamma - q \eqtext{scalar}( \Phi ^{*} U / \gamma )) dt,\notag \\
&= \int ( -m / \gamma,\notag \\
&\quad - q \eqtext{scalar}( (V - \mathbf A \cdot \sigv ) (1 + \mathbf v \cdot \sigv )) ) dt,\notag \\
&= \int ( -m / \gamma - q V + q \mathbf v \cdot \mathbf A ) dt.
\end{align}
\begin{equation}
\mathcal{L} = -m / \gamma - q ( V - \mathbf v \cdot \mathbf A )\qquad \text{Lagrangian}.
\end{equation}
\noindent\textit{Verification.}\par
\begin{align}
\delv{} \mathcal{L} &= m \gamma \mathbf v + q \mathbf A,\notag \\
&= \mathbf P + q \mathbf A.
\end{align}
\begin{align}
d/dt \delv{} \mathcal{L} &= d\mathbf P/dt + q \partial _{t} \mathbf A + q (\mathbf v \cdot \delr{}) \mathbf A, \\
\delr{} \mathcal{L} &= - q \delr{} V + q \delr{} (\mathbf v \cdot \mathbf A).
\end{align}
\noindent\textit{Use the following vector-calculus identity.}\par
\begin{equation}
(\mathbf v(t) \cdot \delr{}) \mathbf A(t, \mathbf r) = \delr{} ( \mathbf v \cdot \mathbf A ) - \mathbf v \times ( \delr{} \times \mathbf A ).
\end{equation}
Solving \(d/dt \delv{} \mathcal{L} = \delr{} \mathcal{L}\) for \(d\mathbf P/dt\) gives the following relation.\par
\begin{align}
d\mathbf P/dt &= - q \delr{} V - q \partial _{t} \mathbf A,\notag \\
&\quad + q ( \delr{} (\mathbf v \cdot \mathbf A) - (\mathbf v \cdot \delr{}) \mathbf A ),\notag \\
&= q ( -\delr{}V - \partial _{t}\mathbf A + \mathbf v \times ( \delr{} \times \mathbf A ) ),\notag \\
&= q ( \mathbf E + \mathbf v \times \mathbf B ),\notag \\
&= m \eqtext{vector}( dU/ds ) / \gamma.
\end{align}
\begin{align}
\mathcal{H} &= (\delv{} \mathcal{L}) \cdot \mathbf v - \mathcal{L},\notag \\
&= m \gamma \mathbf v \cdot \mathbf v + q \mathbf v \cdot \mathbf A + m / \gamma.
\end{align}
\noindent\textit{Adding the potential term contributes } \(qV-q\,\mathbf v\cdot\mathbf A\)\textit{.}\par
\begin{align}
&= m \gamma ( \mathbf v^{2} + 1 / \gamma ^{2} ) + q V,\notag \\
&= m \gamma ( \mathbf v^{2} + 1 - \mathbf v^{2} ) + q V,\notag \\
&= m \gamma + q V.
\end{align}
\noindent\textit{Recall the Hamiltonian property.}\par
\begin{align}
dH/dt &= - \partial _{t} \mathcal{L},\notag \\
d\mathcal{H}/dt &= d(m \gamma )/dt + q dV/dt,\notag \\
&= d(m \gamma )/dt + q \partial _{t}V + q \mathbf v \cdot \delr{} V.
\end{align}
\begin{align}
\partial _{t} H &= q \partial _{t} V - q \mathbf v \cdot \partial _{t} \mathbf A,\notag \\
d(m \gamma )/dt &= q \mathbf v \cdot ( -\delr{} V - \partial _{t} \mathbf A ),\notag \\
&= q \mathbf v \cdot \mathbf E.
\end{align}

\subsubsection*{Electromagnetic Energy-Momentum}
\begin{align}
F &= (\mathbf E + i \mathbf B) \cdot \sigv,\notag \\
F^{\mathrm{R}} &= (\mathbf E - i \mathbf B) \cdot \sigv,\notag \\
T &= \tfrac{1}{2} F F^{\mathrm{R}},\notag \\
&= \tfrac{1}{2} (\mathbf E^{2} + \mathbf B^{2}) + (\mathbf E \times \mathbf B) \cdot \sigv,\notag \\
&= T_{s} + T_{\mathbf v} \cdot \sigv.
\end{align}
\begin{equation}
\partial = \partial _{t} + \sigv \cdot \delr{}.
\end{equation}
\noindent\flag{This is a frame-dependent energy-density/Poynting paravector, not by itself the full electromagnetic stress-energy tensor.}\par
\noindent\textit{Recall for any paravectors Z and Z'.}\par
\begin{align}
\partial (ZZ') &= (\partial Z) Z' + Z (\partial Z') + [\sigv , Z] \cdot \delr{} Z',\notag \\
\partial T &= \tfrac{1}{2} \partial (FF^{\mathrm{R}}),\notag \\
&= \tfrac{1}{2}(\partial F) F^{\mathrm{R}} + \tfrac{1}{2}F (\partial F^{\mathrm{R}}) + \tfrac{1}{2}[\sigv , F] \cdot \delr{} F^{\mathrm{R}}.
\end{align}
\begin{align}
D &= \rho + \mathbf J \cdot \sigv\qquad \text{density paravector}, \\
D^{*} &= \rho - \mathbf J \cdot \sigv, \\
\partial &= \partial _{t} + \sigv \cdot \delr{}, \\
\partial F &= D^{*}, \\
F^{2} &= (\mathbf E^{2} - \mathbf B^{2}) + 2 i (\mathbf E \cdot \mathbf B)\qquad \text{scalar}.
\end{align}

\noindent\textit{Recall for any paravectors Z and Z'.}\par
\begin{align}
\partial (ZZ') &= (\partial Z) Z' + Z (\partial Z') + [\sigv , Z] \cdot \delr{} Z',\notag \\
\partial (F^{2}) &= (\partial F) F + F (\partial F) + [\sigv , F] \cdot \delr{} F,\notag \\
&= D^{*} F + F D^{*} + [\sigv , F] \cdot \delr{} F.
\end{align}
\begin{align}
F &= F_{\mathbf v} \cdot \sigv,\notag \\
[\sigv , F] &= 2 i F_{\mathbf v} \times \sigv,\notag \\
[\sigv , F] \cdot \delr{} F &= -2 i F_{\mathbf v} \cdot (\delr{} \times F_{\mathbf v}),\notag \\
&\quad -2 ( (\delr{} \cdot F_{\mathbf v}) F_{\mathbf v} - \tfrac{1}{2} \delr{} F_{\mathbf v}^{2} ) \cdot \sigv.
\end{align}
\begin{equation}
\square F^{2} = \partial ^{*} \partial (F^{2})\qquad \text{scalar}.
\end{equation}
